\section*{Capítulo \ref{ch:b1:rays-reflection-refraction} --- Óptica geométrica: raios, reflexão, refração}

\begin{solution}{exo:b1:rays-reflection-refraction:1}
$v = c/n$: água \qty{2.26e8}{m/s}, vidro \qty{2.00e8}{m/s}, diamante
\qty{1.24e8}{m/s}. Fibra: $t = nL/c = 1.48 \times 10^3/3.00\times10^8 =
\qty{4.9}{\micro s}$.
\end{solution}

\begin{solution}{exo:b1:rays-reflection-refraction:2}
$\sin i_2 = \sin i_1/1.33$: $30^\circ \to 22.1^\circ$; $60^\circ \to
40.6^\circ$. Da água a $50^\circ$: $1.33\sin 50^\circ = 1.02 > 1$,
nenhum raio refratado --- \omterm{thm:b1:rays-reflection-refraction:tir}{reflexão total} ($i_c = 48.8^\circ$).
\end{solution}

\begin{solution}{exo:b1:rays-reflection-refraction:3}
$i_c = \arcsin(1/n)$: vidro $41.8^\circ$, diamante $24.4^\circ$, água
$48.8^\circ$. O disco de céu abrange $2 \times 48.8^\circ = 97.5^\circ$.
\end{solution}

\begin{solution}{exo:b1:rays-reflection-refraction:4}
Olhos a \qty{1.60}{m}. O raio que vem dos pés chega ao olho após
refletir-se à metade da altura dos olhos, \qty{0.80}{m}; o raio que vem do alto
da cabeça reflete-se a meio caminho entre o topo e os olhos, a \qty{1.65}{m}.
Espelho de \qty{0.80}{m} a \qty{1.65}{m}: altura \qty{0.85}{m}, metade
da altura do corpo --- independente da distância (a regra do ponto médio
vale a qualquer distância).
\end{solution}

\begin{solution}{exo:b1:rays-reflection-refraction:5}
Para ângulos pequenos, a posição aparente é $h' = h/n$ quando se olha do
meio menos refringente, e $h' = nh$ do mais refringente. Moeda: $1.2/1.33 =
\qty{0.90}{m}$. Pescador: $1.33 \times 1.8 = \qty{2.4}{m}$. Justificativa:
para um ponto da superfície à distância horizontal $x$, $\tan i_2 \approx i_2 =
x/h$ e $\tan i_1 \approx i_1 = x/h'$ com $i_1 = n i_2$.
\end{solution}

\begin{solution}{exo:b1:rays-reflection-refraction:6}
$\sin\frac{A + D_{\min}}{2} = 1.52\sin 30^\circ = 0.760$, logo
$\frac{A + D_{\min}}{2} = 49.5^\circ$, $D_{\min} = 38.9^\circ$.
A emergência exige $r' \leq i_c$; além disso $\sin r = \sin i/n \leq 1/n$ dá
$r \leq i_c$; logo $A = r + r' \leq 2i_c$. Aqui $i_c = \arcsin(1/1.52) =
41.1^\circ$, $2i_c = 82.3^\circ > 60^\circ$: tudo bem.
\end{solution}

\begin{solution}{exo:b1:rays-reflection-refraction:7}
$\mathrm{NA} = \sqrt{1.50^2 - 1.48^2} = \sqrt{0.0596} = 0.244$,
$\theta_a = 14.1^\circ$. $\Delta t/L = (1.50/c)(1.50/1.48 - 1) =
5.0\times10^{-9} \times 0.0135 = \qty{68}{ns/km}$. Ao longo de \qty{1}{km}:
$B_{\max} \approx 1/(2 \times \qty{68}{ns}) = \qty{7.4}{Mbit/s}$.
\end{solution}

\begin{solution}{exo:b1:rays-reflection-refraction:8}
$D \approx (n - 1)A$: vermelho $0.510 \times 10^\circ = 5.10^\circ$, azul
$5.30^\circ$; largura do espectro $0.20^\circ$.
\end{solution}

\begin{solution}{exo:b1:rays-reflection-refraction:9}
Entrada: $\sin r = \sin 45^\circ/1.5 = 0.471$, $r = 28.1^\circ$. As duas
faces são paralelas, de modo que a incidência na segunda face vale $r$ e
$\sin i' = n\sin r = \sin i$: $i' = i$, o raio emergente é paralelo ao
incidente. Ao longo do percurso interno de comprimento $e/\cos r$, o
deslocamento perpendicular vale $(e/\cos r)\sin(i - r)$:
$d = 2.0 \times \sin 16.9^\circ/\cos 28.1^\circ = \qty{0.66}{cm}$.
\end{solution}

\begin{solution}{exo:b1:rays-reflection-refraction:10}
Duas refrações ($2(i - r)$) e duas reflexões ($2(\pi - 2r)$):
$D = 2\pi + 2i - 6r$. $\dd D/\dd i = 2 - 6\cos i/(n\cos r) = 0$ dá
$3\cos i = n\cos r$, $9\cos^2 i = n^2 - 1 + \cos^2 i$, $\cos^2 i =
(n^2 - 1)/8$. Com $n = 1.333$: $\cos^2 i = 0.0971$, $i = 71.8^\circ$,
$\sin r = 0.950/1.333$, $r = 45.5^\circ$, $D = 360^\circ + 143.7^\circ -
272.8^\circ = 230.9^\circ$: o arco fica a $230.9^\circ - 180^\circ =
51^\circ$ do ponto anti-solar, fora do primário ($42^\circ$).
Vermelho ($n = 1.331$) a $50.4^\circ$, violeta ($1.343$) a $53.6^\circ$:
as cores estão invertidas, o violeta por fora.
\end{solution}

\begin{solution}{exo:b1:rays-reflection-refraction:11}
$\sin i_c = 1.00026/1.00029 = 0.99997$, $i_c = 89.56^\circ$: os raios
a menos de $0.44^\circ = \qty{7.7}{mrad}$ da superfície da estrada sofrem \omterm{thm:b1:rays-reflection-refraction:tir}{reflexão
total} (em geral, ângulo rasante $\approx \sqrt{2\Delta n/n}$). De
\qty{1.2}{m} de altura, a linha de visada fica rasante nesse ângulo além de
$1.2/0.0077 \approx \qty{160}{m}$: a estrada ali reflete o céu ---
a ``poça''.
\end{solution}

\begin{solution}{exo:b1:rays-reflection-refraction:12}
$t(x) = \dfrac{\sqrt{a^2 + x^2}}{v_1} + \dfrac{\sqrt{b^2 + (d - x)^2}}{v_2}$;
$t'(x) = \dfrac{x}{v_1\sqrt{a^2 + x^2}} - \dfrac{d - x}{v_2\sqrt{b^2 + (d - x)^2}}
= \dfrac{\sin i_1}{v_1} - \dfrac{\sin i_2}{v_2}$, nulo exatamente quando
$\sin i_1/v_1 = \sin i_2/v_2$, isto é,\ $n_1\sin i_1 = n_2\sin i_2$ com
$n = c/v$.
\end{solution}

\begin{solution}{pb:b1:rays-reflection-refraction:1}
\textbf{1.} $n = c/v$; $v_1 = 3.00\times10^8/1.480 = \qty{2.03e8}{m/s}$.

\textbf{2.} $t = Ln_1/c = 5.0\times10^4 \times 1.48/3.00\times10^8 =
\qty{247}{\micro s}$.

\textbf{3.} $\nu = c/\lambda_0 = \qty{1.94e14}{Hz}$, inalterada no vidro;
$\lambda = \lambda_0/n_1 = \qty{1047}{nm}$.

\textbf{4.} Um \omterm{def:b1:rays-reflection-refraction:fiber}{núcleo} nu só guia enquanto a sua superfície estiver
perfeitamente limpa e intocada: poeira, água ou um aperto substituem
localmente o ar por um índice maior e a luz escapa, e os riscos a espalham.
A \omterm{def:b1:rays-reflection-refraction:fiber}{casca} é uma interface selada e controlada, com um degrau de índice escolhido de propósito.

\textbf{5.} $\sin i_c = 1.465/1.480 = 0.9899$, $i_c = 81.8^\circ$.

\textbf{6.} $90^\circ - 81.8^\circ = 8.2^\circ$.

\textbf{7.} $\sin\theta = n_1\sin\theta'$; guiado se e só se $\cos\theta' \geq
n_2/n_1$, isto é,\ $\sin\theta' \leq \sqrt{1 - n_2^2/n_1^2}$, isto é,\
$\sin\theta \leq \sqrt{n_1^2 - n_2^2} = \sqrt{2.1904 - 2.1462} = 0.210$:
$\theta_a = 12.1^\circ$.

\textbf{8.} $\mathrm{NA} = 0.210$. Um degrau maior alarga o cone, de modo que
mais luz de uma fonte divergente (um LED) é captada.

\textbf{9.} O raio é tangente ao círculo de raio $R_b$ (borda
interna) e encontra o círculo de raio $R_b + 2a$ (borda externa); no
triângulo retângulo centro--ponto de tangência--ponto de impacto, o ângulo no
ponto de impacto entre o raio e o raio do círculo (a normal) satisfaz
$\sin i = R_b/(R_b + 2a)$. Guiado se e só se $\sin i \geq n_2/n_1$:
$R_b \geq 2a\,\dfrac{n_2/n_1}{1 - n_2/n_1} = \qty{50}{\micro m} \times 97.6
= \qty{4.9}{mm}$. Curvas mais fechadas vazam.

\textbf{10.} $t_0 = Ln_1/c$; percurso do raio mais inclinado $L/\cos\theta'_{\max} =
Ln_1/n_2$, de modo que $t_1 = Ln_1^2/(n_2 c)$.

\textbf{11.} $\Delta t/L = (n_1/c)(n_1/n_2 - 1) = 4.93\times10^{-9} \times
0.01024 = \qty{50.5}{ns/km}$; ao longo de \qty{2}{km}: \qty{101}{ns}.

\textbf{12.} $B_{\max} \approx 1/(2\Delta t) = 1/\qty{202}{ns} \approx
\qty{5}{Mbit/s}$.

\textbf{13.} $B_{\max} L = 1/(2\,\Delta t/L) = 1/(2 \times
\qty{50.5}{ns/km}) \approx \qty{10}{Mbit.s^{-1}.km}$.

\textbf{14.} Os raios afastados do eixo percorrem caminhos mais longos, mas em
índice menor, portanto mais depressa; com um perfil bem escolhido os dois efeitos se cancelam
e todos os raios chegam juntos.

\textbf{15.} $\qty{9}{\micro m} \approx 6\lambda_0$: o \omterm{def:b1:rays-reflection-refraction:fiber}{núcleo} não é
grande em comparação com o comprimento de onda, de modo que a difração domina e os raios
perdem o sentido. O tratamento ondulatório mostra um único padrão guiado: nenhum
alargamento intermodal; restam apenas a dispersão cromática e a atenuação.

\textbf{16.} \qty{100}{km} $\times$ \qty{0.2}{dB/km} $= \qty{20}{dB}$:
$P/P_0 = 10^{-2}$, um por cento.

\textbf{17.} $10\log(P_0/P_{\min}) = \qty{30}{dB}$, logo $L_{\max} =
30/0.20 = \qty{150}{km}$.

\textbf{18.} \qty{100}{km} custam \qty{20}{dB}: margem de \qty{10}{dB}, um
fator $10$ em potência.

\textbf{19.} LED: $\Delta t_c = 17 \times 40 \times 100 = \qty{68000}{ps}
= \qty{68}{ns}$, $B \approx 1/\qty{136}{ns} = \qty{7}{Mbit/s}$. Laser:
$\Delta t_c = \qty{170}{ps}$, $B \approx 1/\qty{340}{ps} =
\qty{2.9}{Gbit/s}$.

\textbf{20.} $30/2.0 = \qty{15}{km}$, dez vezes menos. A sílica é mais
transparente perto de \qty{1550}{nm} (o seu mínimo de atenuação), e ali existem
amplificadores ópticos.

\textbf{21.} $30/30 = \qty{1}{km}$ entre amplificadores, contra
\qty{150}{km}: a fibra vence por um fator $150$.

\textbf{22.} Intermodal: \qty{101}{ns} ($\to$ \qty{5}{Mbit/s});
cromática (laser): $17 \times 0.1 \times 2 = \qty{3.4}{ps}$;
atenuação \qty{0.4}{dB}, desprezível. O alargamento intermodal limita:
\qty{5}{Mbit/s}.

\textbf{23.} Nenhum alargamento intermodal; cromático \qty{170}{ps} $\to$
\qty{2.9}{Gbit/s}; atenuação \qty{20}{dB}, dentro do orçamento. A dispersão
cromática limita, em cerca de \qty{3}{Gbit/s}.

\textbf{24.} $80 \times 2.9 \approx \qty{230}{Gbit/s}$.

\textbf{25.} Índice degrau: \qty{10}{Mbit.s^{-1}.km}; canal monomodo:
$2.9 \times 100 = \qty{290}{Gbit.s^{-1}.km}$ --- um fator de cerca de
$3\times10^4$. Um fato: um único caminho guiado, portanto nenhum alargamento
geométrico (intermodal) dos tempos de chegada.
\end{solution}
